\section{Optimized finite backgrounds below cutoff one}\label{sec:background}
Fix $0<\rho<1$ and, throughout this section, set
\[
 s=\frac{1+\rho}{1-\rho},\qquad k_s=\frac{s-1}{s},\qquad
 C_s=\frac{s^s}{(s-1)^{s-1}}.
\]
The cutoff-one result does not classify a fixed cutoff $K<1$ with
a tail budget greater than one. We now combine long parities with the
balanced critical tilt to give a sub-two collapse region. The sharp
optimization below concerns a scalar background law, not the optimal
spectral boundary among all densities.

\begin{proposition}[Scalar background tail optimum]\label{prop:background-optimum}
For $s>1$ and $0<K<1$, among nonnegative random variables $A$ with
$\mathbb E A=1$,
\[
 \inf\mathbb E[A^s\mathbf1_{A>K}]=m_s(K),\qquad
 m_s(K)=\begin{cases}
 1,&0<K\le k_s,\\
 C_s(1-K)K^{s-1},&k_s<K<1.
 \end{cases}
\]
For $K\le k_s$ the constant law $A=1$ attains the infimum. For $K>k_s$
an attaining law has values $K$ and $v_*=sK/(s-1)$, with high probability
$p_*=(s-1)(1-K)/K$. Replacing the strict event by $A\ge K$ leaves
the infimum unchanged, but that infimum is not attained when $K>k_s$.
\end{proposition}
\begin{proof}
For $K\le k_s$, convexity gives $y^s\ge1+s(y-1)$. If $y\le K$, the
right side is nonpositive. Hence $y^s\mathbf1_{y>K}\ge1+s(y-1)$
everywhere. Taking expectations gives the lower bound one, attained
by $A=1$. The affine minorant is integrable even if the left side has
infinite expectation.

For $K>k_s$, the derivative of $y^s/(y-K)$ on $y>K$ has the sign
of $(s-1)y-sK$. Its unique minimum is at $v_*$ and its value is
$c=C_sK^{s-1}$. Thus
\[
 y^s\mathbf1_{y>K}\ge c(y-K)\qquad(y\ge0).
\]
Integration gives $c(1-K)$. The stated probabilities lie in $(0,1)$,
have mean $K+p_*(v_*-K)=1$, and put mass only at the two equality
points of the minorant. At $K=k_s$ this law degenerates to $A=1$.

For the non-strict event the same lower bound holds. When $K>k_s$,
choose $u<K$ tending to $K$ with $u>k_s$, high value $su/(s-1)$,
and high probability $(s-1)(1-u)/u$. The high value exceeds $K$
eventually, and the cost tends to $C_s(1-K)K^{s-1}$. Thus the
infimum is unchanged. Equality in the affine minorant for the
non-strict integrand is possible only at $v_*$: below $K$ its
right side is strictly negative, and at $K$ its left side is $K^s>0$.
Concentration at $v_*>1$ cannot have mean one, proving nonattainment.
For $K\le k_s$, $A=1$ still attains one.
\end{proof}

This is an elementary instance of classical moment optimization by
pointwise affine bounds; see Bertsimas and Popescu~\cite{BertsimasPopescu}
for the general moment-duality setting. Neither that method nor the
formula's external priority is claimed as new here.

\begin{theorem}[Sub-two tail budgets with finite parity backgrounds]
\label{thm:background-collapse}
Fix $k_s<K<1$, $0<b<K$, and
\[
 1+C_s(1-K)K^{s-1}<\tau<2.
\]
There are strictly positive densities $r_L$ on every Boolean cube
of integer dimension $L\ge1$ satisfying
\[
 \mathbb E r_L=1,\qquad r_L\ge b,\qquad
 \mathbb E[r_L^s\mathbf1_{r_L>K}]\le\tau,
\]
such that both leading eigenvalues of $T_{r_L}$ tend to one and
$\gap(r_L)\to0$. The densities use global interactions. No marginal
matching, locality, equality-budget attainability or optimal universal
spectral threshold is asserted.
\end{theorem}
\begin{proof}
\emph{A fixed dyadic background.}
Write $g(u)=C_s(1-u)u^{s-1}$. On $(k_s,1)$,
\[
 g'(u)=C_su^{s-2}((s-1)-su)<0,
\]
and its endpoint values are one and zero. Strict budget slack permits
a fixed $\max(b,k_s)<u<K$ with $1+g(u)<\tau$. The optimal high
probability at lower value $u$ is $p_u=(s-1)(1-u)/u$.
For $p\in(0,1)$ define $v(p)=u+(1-u)/p$. The mean
$(1-p)u+pv(p)$ is exactly one and $pv(p)^s$ is continuous in $p$.
Choose a fixed dyadic $p=j/2^k$ sufficiently close to $p_u$ that
\[
 1+pv^s<\tau,\qquad v=v(p)>1>K,
 \qquad k\ge1,\quad 1\le j<2^k.
\]
On $k$ unbiased signs choose a function $F$ equal to $v$ at $j$
patterns and equal to $u$ elsewhere. Then $\mathbb EF=1$.
Replace each sign by the parity of a disjoint block of $M$ unbiased
bits, obtaining $a_M$ on $kM$ bits. Those parities are independent
unbiased signs, so $a_M$ has exactly the same two-point law.
All $u,v,p,k$ are fixed before the dimension grows.

\emph{An operator-norm estimate for every character.}
Let $Q_d=P_{\sqrt\rho,d}$ and let $\chi_U$ be a Walsh character.
For any Walsh basis vector $\chi_S$,
\[
 Q_dM_{\chi_U}Q_d\chi_S
 =\rho^{(|S|+|S\mathbin\triangle U|)/2}\chi_{S\mathbin\triangle U}
 =\rho^{|U|/2+|S\setminus U|}\chi_{S\mathbin\triangle U}.
\]
This is a weighted permutation of an orthonormal basis. Its norm is
exactly $\rho^{|U|/2}$, attained for $S\subset U$.
Writing the finite Walsh expansion
$F=1+\sum_{I\ne\varnothing}\widehat F(I)\chi_I$ and putting
$D_F=\sum_{I\ne\varnothing}|\widehat F(I)|$, block substitution gives
\[
 \|Q_{kM}M_{a_M}Q_{kM}-P_{\rho,kM}\|
 \le \sum_{I\ne\varnothing}|\widehat F(I)|\rho^{M|I|/2}
 \le D_F\rho^{M/2}\longrightarrow0.                 \tag{31}
\]
Large fixed amplitudes or coefficients do not obstruct this limit.

\emph{The inherited critical family and actual eigenvalues.}
Use the balanced critical family in Theorem~\ref{thm:critical-exact}:
with fixed small $d>0$, $\delta_n=dn^{-1/3}$ and
\[
 N(\delta)=\left(\frac{(1+\delta)^s+(1-\delta)^s}{2}\right)^{1/s},
 \qquad R_n(y)=\prod_{i=1}^n\frac{1+\delta_n y_i}{N(\delta_n)},
\]
put
\[
 \kappa(\delta)=\frac{1+\rho+
 \sqrt{(1-\rho)^2+4\rho\delta^2}}{2N(\delta)},\quad
 \beta_n=\kappa(\delta_n)^{-n},\quad
 Z_n=1+\beta_n\mathbb ER_n,\quad
 w_n=\frac{1+\beta_nR_n}{Z_n}.
\]
The established balanced Gram estimates give the actual limits
$\lambda_1(T_{w_n}),\lambda_2(T_{w_n})\to1$, not merely two
trial-vector bounds. Also $\beta_n,Z_n\to1$, $\mathbb ER_n^s=1$,
and $\mathbb ER_n^t\to0$ for every fixed $0<t<s$.
The latter follows from the one-site expansion
\[
 \log\mathbb E\left[\frac{1+\delta Y}{N(\delta)}\right]^t
 =\frac{t(t-s)}2\delta^2+O(\delta^4),
\]
since $n\delta_n^2\to\infty$ and $n\delta_n^4\to0$.

For sufficiently large $L$ take $M=\lfloor L/(2k)\rfloor$ and
$n=L-kM$. Both tend to infinity. On the two independent blocks set
\[
 r_L(x,y)=\frac{a_M(x)+\beta_nR_n(y)}{Z_n}.
\]
Its mean is one and its minimum is at least $u/Z_n$, eventually
greater than $b$. The matrices $S_r=Q_LM_rQ_L$ and $T_r$ have the
same full spectrum by the $AA^*$ and $A^*A$ identity with
$A=M_{\sqrt r}Q_L$. Linearity and tensorization show
\[
 \left\|S_{r_L}-P_{\rho,kM}\otimes S_{w_n}\right\|
 \le D_F\rho^{M/2}/Z_n\longrightarrow0.
\]
The tensor comparison has largest eigenvalue $\lambda_1(S_{w_n})$
and second eigenvalue between $\lambda_2(S_{w_n})$ and
$\lambda_1(S_{w_n})$: use the constant background eigenvector
and the first two tilt eigenvectors for the lower bound. Positivity
and the fact that the top background eigenvalue is one give the
upper bound. Both leading eigenvalues therefore tend to one.
The self-adjoint eigenvalue perturbation inequality transfers both
limits to $S_{r_L}$ and hence $T_{r_L}$. Each finite matrix remains
strictly positive, with a simple Perron root.

\emph{Noninteger moments and the moving strict cutoff.}
Let $\theta=\min(1,s/2)$. For $c>0,z\ge0$,
\[
 0\le(c+z)^s-c^s-z^s
 \le s2^{s-1}\big(c^{s-\theta}z^\theta+
                   c^\theta z^{s-\theta}\big).       \tag{32}
\]
To prove this, divide by $c^s$ and put $y=z/c$. For $y\le1$,
the mean value theorem bounds $(1+y)^s-1-y^s$ by
$s2^{s-1}y\le s2^{s-1}y^\theta$. For $y\ge1$ interchange
$c$ and $z$. Nonnegativity follows from convexity. Both powers
$\theta,s-\theta$ lie strictly between zero and $s$.
For $z=\beta_nR_n$, the mixed expectations vanish and
$\mathbb Ez^s\to1$. Thus $\mathbb E(c+\beta_nR_n)^s\to c^s+1$.

In the high sector, $v/Z_n>K$ eventually, so the entire sector
contributes and its conditional tail moment tends to $v^s+1$.
In the low sector the event is exactly $R_n>c_n$, where
$c_n=(KZ_n-u)/\beta_n\to K-u>0$. Fix positive finite
$c_-,c_+$ containing these thresholds eventually. For $0<t<s$,
\[
 \mathbb P(R_n>c_n)\le c_-^{-t}\mathbb ER_n^t\to0,
 \qquad
 \mathbb E[R_n^s\mathbf1_{R_n\le c_n}]
 \le c_+^{s-t}\mathbb ER_n^t\to0.
\]
Consequently the strict event retains asymptotic $s$-moment one.
In (32) restricted to this event, the constant term vanishes with
its probability, and the mixed terms are bounded by their vanishing
unrestricted expectations. After normalization, the low-sector
tail moment tends to one. Averaging the fixed sector probabilities,
\[
 \mathbb E[r_L^s\mathbf1_{r_L>K}]
 \longrightarrow (1-p)+p(v^s+1)=1+pv^s<\tau.
\]
This calculation does not infer moment convergence from convergence
in probability. Indeed the full moment tends to
$1+(1-p)u^s+pv^s>2$, by strict Jensen for the nonconstant background.

Choose a common $L_0$ beyond which both blocks are positive, the
floor holds and the tail cost is at most $\tau$. For $1\le L<L_0$
set $r_L=1$. Its tail cost is one and its floor is valid since
$b<K<1$. This finite replacement preserves all limits and proves
the assertion in every dimension. No effective value of $L_0$ is claimed.
\end{proof}

\subsection*{Comparison and the unresolved spectral interval}
The half-parity background with low value approaching $K$ and high
value approaching $2-K$ has limiting scalar cost $(2-K)^s/2$.
It gives collapse whenever
\[
 2-2^{1/s}<K<1,\qquad 0<b<K,\qquad
 1+(2-K)^s/2<\tau<2.
\]
This is a specialization of the same proof with $p=1/2$ and
$u<K$ chosen with strict slack. Proposition~\ref{prop:background-optimum}
shows that the optimized cost is no larger. Equality occurs only
when $2-K=sK/(s-1)$, that is $K=2(s-1)/(2s-1)$; outside this
isolated value the cost improvement is strict. Strict convexity
of $2^t$ below its chord on $0<t<1$ gives
$2^{1/s}<1+1/s$. Hence the new interval $k_s<K<1$ strictly
extends the half-parity sub-two interval.

On the other hand, any mean-one density with tail budget at most
$\tau$ obeys $\mathbb Er^s\le K^{s-1}+\tau$, since
$r^s\mathbf1_{r\le K}\le K^{s-1}r\mathbf1_{r\le K}$.
The established total-moment estimate therefore gives
\[
 \gap(r)\ge1-\left(\frac{K^{s-1}+\tau}{2}\right)^{1/s}>0
 \quad\hbox{if }\tau<2-K^{s-1}.
\]
Neither this sufficient boundary nor $1+m_s(K)$ is proved to be
the universal spectral threshold. In particular, adding an escaped
unit to the scalar optimum is justified for the construction above;
it is not a lower bound for every possible collapsing family.
The equality budget, the intervening interval and the floor $b=K$
remain unclassified here. For $K<b<1$ the event contains every point
and reduces to the total-moment condition. At $b=1$ mean one forces
$r=1$ and gap $1-\rho$; at $b>1$ the class is empty. The earlier
source309A omitted $b<1$ in its additional comparison; its failed
gate and the corrected309B report remain preserved, not relabelled.
